\setlength{\baselineskip}{20pt}

\begin{englishabstract}
The switch rails, stock rails, and frogs of railway turnouts are critical components prone to surface and near-surface damage during long-term service, posing a threat to operational safety. Existing inspection methods face limitations such as difficult deployment and insufficient anti-interference capability in the spatially constrained turnout environment. Electromagnetic Tomography (EMT) is a non-contact, non-destructive testing technique based on alternating electromagnetic field induction, featuring planar-array configurability and suitability for near-surface defect detection, and thus holds application potential for turnout damage inspection. However, turnout steel is a highly conductive and ferromagnetic target, where defect-induced electromagnetic perturbations are extremely weak and the EMT inverse problem is severely ill-posed and nonlinear. Consequently, conventional sensitivity-matrix-based linear reconstruction methods struggle to simultaneously achieve satisfactory accuracy, robustness, and engineering applicability.

This thesis presents an investigation into the design of a portable EMT hardware platform and intelligent imaging algorithms, establishing a complete validation pipeline from electromagnetic response acquisition to damage imaging display. On the algorithmic side, to address the insufficient adaptation of the classical Deep Operator Network (DeepONet) Branch--Trunk structure to EMT multi-excitation structured inputs and continuous-domain outputs, a Shared-Conditional-Feature Deep Operator Network (SCF-DeepONet) is proposed. The network maps multi-view complex differential voltage responses into global conditional features, excitation-specific features, and a two-dimensional local feature map through a shared observation encoder. A primary branch for two-dimensional defect heatmap prediction and an auxiliary branch for 3d magnetic vector potential difference-field reconstruction are constructed in parallel. Ablation experiments demonstrate that, under identical training conditions, SCF-DeepONet achieves a mean Intersection over Union (IoU) of 81.53\% on the independent test set, representing a 13.73\% absolute improvement over the Vanilla DeepONet baseline (67.79\%). The Worst-10 Avg IoU improves from 37.09\% to 51.34\%, confirming the effectiveness of the proposed architecture for ferromagnetic EMT damage detection. Furthermore, three-dimensional difference-field auxiliary supervision is introduced as a training-stage physics-aware regularization and a progressive pre-experiment for Physics-Informed Neural Networks (PINN): a four-stage curriculum training strategy—comprising heatmap pretraining, heatmap-field joint training, heatmap refinement, and Partial Differential Equation (PDE)-guided weak physical fine-tuning—is designed. Experimental results indicate that auxiliary supervision improves the mean IoU on zero-shot circular defect samples from 34.05\% to 40.31\%, enhancing out-of-distribution generalization without necessarily increasing the average accuracy on in-distribution rectangular samples, and also providing empirical foundation and initialization prerequisite for subsequent strict PDE constraints.

On the hardware side, a portable EMT detection platform is designed and implemented. The system adopts a dual-board architecture consisting of a measurement-and-control board and a signal-switching board, featuring a $4 \times 4$ planar coil array sensor, a heterogeneous transmit-receive high-current multiplexing network, an excitation generation and dual-stage power amplification chain based on Direct Digital Synthesis (DDS), a dual-phase lock-in demodulation module and a high-precision analog-to-digital conversion module, together with main-controller communication, power management, and Printed Circuit Board (PCB) anti-interference design. The system can stably acquire multi-channel complex mutual inductance responses at the 100 kHz operating frequency, with a complete frame acquisition period of approximately 1.46 s. Hardware performance tests confirm stable excitation signals, reliable lock-in demodulation, and good channel consistency.

	
\noindent\englishkeywords{electromagnetic tomography; non-destructive testing; deep operator network; multi-task learning; physics-informed auxiliary learning}
\end{englishabstract}